\subsection{PX4 Architecture and Software Fault Localization}

PX4 adopts a modular flight-control architecture in which estimators, controllers, sensor drivers, and other software components exchange data through the uORB publish--subscribe middleware \cite{meier2015px4,px4uorbdocs2026}. Many of these uORB topics are also recorded in ULog flight telemetry, providing an explicit relationship between runtime flight data and the software modules that produce or consume the corresponding messages. This architecture makes uORB a useful interface for relating telemetry-level observations to the internal structure of the flight-control software.

Software fault localization aims to reduce debugging effort by ranking program entities according to their likelihood of being associated with a failure. Traditional approaches commonly exploit execution spectra, test outcomes, or other program-level evidence to rank suspicious statements, methods, or files \cite{wong2016survey,jones2002visualization}. Subsequent studies have also investigated the use of seeded faults and real defects to evaluate localization effectiveness under controlled ground truth \cite{pearson2017evaluating,papadakis2019mutation}. These techniques provide useful principles for narrowing the software inspection space, but they generally assume direct software-level execution or testing evidence.

In telemetry-driven UAV diagnosis, such software-level evidence is usually unavailable during real flight. Instead, the available observations consist of abnormal telemetry signals and model-derived diagnostic evidence. The correspondence between uORB topics and PX4 modules therefore provides a natural bridge between telemetry analysis and software inspection. However, this connection has received limited attention in existing UAV anomaly-diagnosis studies. HS-AeroTS builds on this architectural relationship by linking topic-level telemetry evidence with firmware-specific PX4 module relations to support software-module inspection.
